% Discussion. Short taxonomy reprise; limitations as scientific scope;
% open problems as an enumerate, not a sixteen-row table.

\section{Discussion}
\label{var:sec:discussion}

\subsection{What the collection shows}
\label{var:sec:pattern}

The sorting of \cref{var:fig:taxonomy} is the argument, and it is worth being
exact about the last cell. A rate that moves because the process moved costs
the generating function, and with it most of what the preceding chapters were
able to do --- unless what the process acts on is a single random time rather
than a moving rate. Between those two poles, the interesting object in the
collection is the Pimentel chain: a process that repeatedly restores its own
criticality is not the same thing as one that approaches an equilibrium, and
\cref{var:fig:pim-critical} is the picture of that difference.

\Cref{var:sec:twophase} sits where it does because the phase ends at the
rupture of the host cell, which the process itself brings on at rate
$\delta_1 X_t$. That is feedback of the same kind as \cref{var:sec:rrrr}'s.
What differs is where the feedback is allowed to show up. In every other
model the state bends a rate, and the generating function stops closing.
Here it sets a single random \emph{time}, and the burst-size law has already
integrated that time out --- what crosses the join is a distribution over
$\Kb$ that has forgotten when the rupture came.

That is the trade the chapter makes visible, and it is sharper than feedback
against no feedback. A process may act on its own rates as much as it likes
and still be solvable, provided what survives the interaction is a random
time or a random count rather than a moving rate. It is continuous bending,
not feedback as such, that costs the generating function.

\subsection{Limitations}
\label{var:sec:limitations}

Five of the seven models have never been compared against data, and for three
of those it is not obvious what data would serve. The Pimentel+ specification
of \cref{var:sec:pimentelplus} has not been simulated, so the behaviour
attributed to it is what the construction was intended to produce rather than
anything observed of it. The public-good threshold $K^\ast$ of
\cref{var:fig:kstar} is a numerical estimate on one parameter slice, not a
comparison against data.

The ecological models of \cref{var:sec:speciation,var:sec:rrrr} use
``species'' as a unit whose delimitation is contested in the literature the
speciation section responds to. \Cref{var:sec:rrrr} is explicit that its
$X_t$ might equally count individuals, species or organisational units.

None of \cref{var:sec:pimentelplus,var:sec:rrrr,var:sec:plague,%
var:sec:publicgood} carries an existence result for the behaviour its
motivation suggests. That the vitality coupling gives a sustained oscillation
rather than one collapse, and that the fitness penalty of
\cref{var:eq:fitpenalty} can reverse a dominance, remain conjectures
supported by the construction.

The organism-level claims of \cref{var:sec:publicgood,var:sec:twophase} are
hypotheses about \yp{} compatible with the cited intracellular literature,
not inferences this chapter can check.

\subsection{Open problems}
\label{var:sec:openproblems}

\begin{enumerate}[leftmargin=1.6em,itemsep=0.35em]
\item \textbf{The mixture at $\beta_5$.} The law of $\Bint(t)$ under a
randomly switching capacity, and the resulting mixture of geometrics
\cref{var:eq:mixture}. Tractable: a Markov-modulated integral.
\item \textbf{Identifiability.} What a time series of $S_t$ identifies,
given that only $\Bint$ enters the law, and whether $\beta_1$ and $\beta_3$
can be told apart from one.
\item \textbf{Whether the Pimentel oscillation persists.} Every realisation
run here eventually fixated. Whether the expected number of reversals before
fixation grows without bound in $\alpha$, $\varepsilon$ or $N$, or saturates,
is the question the model was built to answer.
\item \textbf{A diffusion limit} in which \cref{var:eq:icrit} is a moving
unstable equilibrium, so that the divergence of nearby trajectories becomes
a rate.
\item \textbf{The scaling of $K^\ast$.} \Cref{var:fig:kstar} is one slice.
How the threshold moves with $\alpha$ and $c/g$ is open.
\item \textbf{Against the threshold budget.} How $\ex{Z}$ here compares
against $\vartheta^{\ast}$ of \cref{path:sec:ypestis}: whether the $\delta_1$
that puts $\ex{Z}$ above one is the $\delta_1$ that puts the switch where
the removal budget puts it.
\end{enumerate}

Of these, the mixture and the $\vartheta^{\ast}$ comparison would repay
attention first: each is well posed, and each closes a section that currently
stops one step short of an answer.

\subsection{Connections}
\label{var:sec:connections}

\Cref{path:sec:forward} identifies a selection gradient between release
phenotypes and hands the study of variability itself to this chapter.
\Cref{var:sec:speciation,var:sec:pimentel,var:sec:pimentelplus} are what
discharges that: variability as the object rather than as a property of a
fixed mechanism.

Two of the models point back into the core theory. The quasi-stationary
question of \cref{var:sec:rrrr} is \cref{dist:thm:qsd}'s own question asked
of a state-dependent catastrophe rate; the extinction probability of
\cref{var:sec:plague} wants the first-step analysis
\cref{path:sec:rsfirststep} runs on a two-dimensional lattice, applied to a
lattice of a different shape.

The two \yp{} models point forward. \Cref{var:sec:publicgood} offers a
mechanism for the intracellular advantage that \ChPath{} infers but does not
explain, and \cref{var:sec:twophase} offers $\ex{Z}$ for the consequence of
the two-phase strategy at the level of an epidemic. The sixth open problem
above is what would join them to it.

\FloatBarrier
